% table_ablation.tex  — Table I: Quantization method comparison
% Updated 2026-05-24 with corrected FP32 baseline (PPL=9.35, not 6.92) and
% WikiText-2 test-set column (non-overlapping windows, stride=256, 325,693 tokens).
%
% Usage: % table_ablation.tex  — Table I: Quantization method comparison
% Updated 2026-05-24 with corrected FP32 baseline (PPL=9.35, not 6.92) and
% WikiText-2 test-set column (non-overlapping windows, stride=256, 325,693 tokens).
%
% Usage: % table_ablation.tex  — Table I: Quantization method comparison
% Updated 2026-05-24 with corrected FP32 baseline (PPL=9.35, not 6.92) and
% WikiText-2 test-set column (non-overlapping windows, stride=256, 325,693 tokens).
%
% Usage: % table_ablation.tex  — Table I: Quantization method comparison
% Updated 2026-05-24 with corrected FP32 baseline (PPL=9.35, not 6.92) and
% WikiText-2 test-set column (non-overlapping windows, stride=256, 325,693 tokens).
%
% Usage: \input{tables/table_ablation}

\begin{table*}[t]
\centering
\caption{Quantization method comparison on TinyStories-15M.
         All CTLE/INT4 methods use 4-bit weight storage.
         CTLE uses $k{=}16$ codebook entries per tensor;
         INT4U uses one global scale per tensor;
         INT4BW uses one scale per group of $G{=}32$ columns
         ($\approx$5\,bits/weight effective).
         NLL and PPL\,(TS) are evaluated on our 10-sentence TinyStories corpus
         (192 tokens scored).
         PPL\,(WT2) is evaluated on the WikiText-2 test split
         (325,693 tokens, non-overlapping windows of 256 tokens, stride\,=\,256).
         Both corpora use the LLaMA SentencePiece tokenizer.
         \textbf{Important:} TinyStories-15M was trained on children-story text,
         not Wikipedia; WikiText-2 PPL values reflect domain mismatch as well as
         compression error and are \emph{not} directly comparable to numbers
         reported for general-purpose LLMs (OPT, LLaMA).
         $\dagger$~PSO/GA use K-means warm start.}
\label{tab:ablation}
\setlength{\tabcolsep}{5pt}
\renewcommand{\arraystretch}{1.15}
\begin{tabular}{lcccccccc}
\toprule
\textbf{Method} & \textbf{Bits/w} & \textbf{Size (MB)} & \textbf{Ratio} &
\textbf{NLL\,$\downarrow$} & \textbf{PPL (TS)\,$\downarrow$} &
\textbf{PPL (WT2)\,$\downarrow$} & \textbf{Time (s)} \\
\midrule
FP32 (baseline)     & 32   & 60.77 & $1.00\times$  & 2.236 & \phantom{00}9.35  & \phantom{0}6{,}127  & —     \\
\midrule
\multicolumn{8}{l}{\textit{Standard INT4 quantization (baselines)}} \\
INT4 Uniform        &  4   &  7.61 & $7.99\times$  & 9.312 & 11{,}070          & 533{,}069           & \phantom{0}${<}1$ \\
INT4 Block-wise     & $\approx$5 & 9.51 & $6.39\times$ & 2.371 & \textbf{10.71} & \textbf{\phantom{0}7{,}120} & \phantom{0}${<}1$ \\
\midrule
\multicolumn{8}{l}{\textit{CTLE — learned codebook (this work)}} \\
CTLE + K-means      &  4   &  7.61 & $7.98\times$  & 2.957 & 19.23           & 19{,}239            & \phantom{0}73.6 \\
CTLE + GA$^\dagger$     &  4   &  7.61 & $7.98\times$  & 2.944 & 18.98           & 19{,}114            & 207.7 \\
CTLE + PSO$^\dagger$    &  4   &  7.61 & $7.98\times$  & 2.768 & 15.93           & 16{,}827            & 198.3 \\
\bottomrule
\end{tabular}
%
\smallskip
\begin{minipage}{\textwidth}
  \footnotesize
  \textbf{INT4U collapses} (PPL\,$>$11\,000 / 533\,000) because a single global scale
  cannot represent the heterogeneous distribution of the 9.2\,M-entry embedding matrix.
  \textbf{INT4BW} recovers near-FP32 quality
  (TS: $+14.5\,\%$; WT2: $+16.2\,\%$ vs.\ FP32) at the cost of 25\,\% larger footprint
  (9.51\,MB vs.\ 7.61\,MB) and a 31\,\% throughput penalty on RISC-V
  (see Table~\ref{tab:hw}).
  \textbf{CTLE-PSO} achieves competitive PPL at identical 4-bit storage and
  streams weights directly from Flash with a fixed 64-byte LUT per tensor,
  eliminating per-group scale-table overhead at inference time.
  PSO reduces PPL by \textbf{17.2\,\%} vs.\ K-means on TinyStories and
  \textbf{12.5\,\%} on WikiText-2, confirming the improvement is
  corpus-independent.
\end{minipage}
\end{table*}


\begin{table*}[t]
\centering
\caption{Quantization method comparison on TinyStories-15M.
         All CTLE/INT4 methods use 4-bit weight storage.
         CTLE uses $k{=}16$ codebook entries per tensor;
         INT4U uses one global scale per tensor;
         INT4BW uses one scale per group of $G{=}32$ columns
         ($\approx$5\,bits/weight effective).
         NLL and PPL\,(TS) are evaluated on our 10-sentence TinyStories corpus
         (192 tokens scored).
         PPL\,(WT2) is evaluated on the WikiText-2 test split
         (325,693 tokens, non-overlapping windows of 256 tokens, stride\,=\,256).
         Both corpora use the LLaMA SentencePiece tokenizer.
         \textbf{Important:} TinyStories-15M was trained on children-story text,
         not Wikipedia; WikiText-2 PPL values reflect domain mismatch as well as
         compression error and are \emph{not} directly comparable to numbers
         reported for general-purpose LLMs (OPT, LLaMA).
         $\dagger$~PSO/GA use K-means warm start.}
\label{tab:ablation}
\setlength{\tabcolsep}{5pt}
\renewcommand{\arraystretch}{1.15}
\begin{tabular}{lcccccccc}
\toprule
\textbf{Method} & \textbf{Bits/w} & \textbf{Size (MB)} & \textbf{Ratio} &
\textbf{NLL\,$\downarrow$} & \textbf{PPL (TS)\,$\downarrow$} &
\textbf{PPL (WT2)\,$\downarrow$} & \textbf{Time (s)} \\
\midrule
FP32 (baseline)     & 32   & 60.77 & $1.00\times$  & 2.236 & \phantom{00}9.35  & \phantom{0}6{,}127  & —     \\
\midrule
\multicolumn{8}{l}{\textit{Standard INT4 quantization (baselines)}} \\
INT4 Uniform        &  4   &  7.61 & $7.99\times$  & 9.312 & 11{,}070          & 533{,}069           & \phantom{0}${<}1$ \\
INT4 Block-wise     & $\approx$5 & 9.51 & $6.39\times$ & 2.371 & \textbf{10.71} & \textbf{\phantom{0}7{,}120} & \phantom{0}${<}1$ \\
\midrule
\multicolumn{8}{l}{\textit{CTLE — learned codebook (this work)}} \\
CTLE + K-means      &  4   &  7.61 & $7.98\times$  & 2.957 & 19.23           & 19{,}239            & \phantom{0}73.6 \\
CTLE + GA$^\dagger$     &  4   &  7.61 & $7.98\times$  & 2.944 & 18.98           & 19{,}114            & 207.7 \\
CTLE + PSO$^\dagger$    &  4   &  7.61 & $7.98\times$  & 2.768 & 15.93           & 16{,}827            & 198.3 \\
\bottomrule
\end{tabular}
%
\smallskip
\begin{minipage}{\textwidth}
  \footnotesize
  \textbf{INT4U collapses} (PPL\,$>$11\,000 / 533\,000) because a single global scale
  cannot represent the heterogeneous distribution of the 9.2\,M-entry embedding matrix.
  \textbf{INT4BW} recovers near-FP32 quality
  (TS: $+14.5\,\%$; WT2: $+16.2\,\%$ vs.\ FP32) at the cost of 25\,\% larger footprint
  (9.51\,MB vs.\ 7.61\,MB) and a 31\,\% throughput penalty on RISC-V
  (see Table~\ref{tab:hw}).
  \textbf{CTLE-PSO} achieves competitive PPL at identical 4-bit storage and
  streams weights directly from Flash with a fixed 64-byte LUT per tensor,
  eliminating per-group scale-table overhead at inference time.
  PSO reduces PPL by \textbf{17.2\,\%} vs.\ K-means on TinyStories and
  \textbf{12.5\,\%} on WikiText-2, confirming the improvement is
  corpus-independent.
\end{minipage}
\end{table*}


\begin{table*}[t]
\centering
\caption{Quantization method comparison on TinyStories-15M.
         All CTLE/INT4 methods use 4-bit weight storage.
         CTLE uses $k{=}16$ codebook entries per tensor;
         INT4U uses one global scale per tensor;
         INT4BW uses one scale per group of $G{=}32$ columns
         ($\approx$5\,bits/weight effective).
         NLL and PPL\,(TS) are evaluated on our 10-sentence TinyStories corpus
         (192 tokens scored).
         PPL\,(WT2) is evaluated on the WikiText-2 test split
         (325,693 tokens, non-overlapping windows of 256 tokens, stride\,=\,256).
         Both corpora use the LLaMA SentencePiece tokenizer.
         \textbf{Important:} TinyStories-15M was trained on children-story text,
         not Wikipedia; WikiText-2 PPL values reflect domain mismatch as well as
         compression error and are \emph{not} directly comparable to numbers
         reported for general-purpose LLMs (OPT, LLaMA).
         $\dagger$~PSO/GA use K-means warm start.}
\label{tab:ablation}
\setlength{\tabcolsep}{5pt}
\renewcommand{\arraystretch}{1.15}
\begin{tabular}{lcccccccc}
\toprule
\textbf{Method} & \textbf{Bits/w} & \textbf{Size (MB)} & \textbf{Ratio} &
\textbf{NLL\,$\downarrow$} & \textbf{PPL (TS)\,$\downarrow$} &
\textbf{PPL (WT2)\,$\downarrow$} & \textbf{Time (s)} \\
\midrule
FP32 (baseline)     & 32   & 60.77 & $1.00\times$  & 2.236 & \phantom{00}9.35  & \phantom{0}6{,}127  & —     \\
\midrule
\multicolumn{8}{l}{\textit{Standard INT4 quantization (baselines)}} \\
INT4 Uniform        &  4   &  7.61 & $7.99\times$  & 9.312 & 11{,}070          & 533{,}069           & \phantom{0}${<}1$ \\
INT4 Block-wise     & $\approx$5 & 9.51 & $6.39\times$ & 2.371 & \textbf{10.71} & \textbf{\phantom{0}7{,}120} & \phantom{0}${<}1$ \\
\midrule
\multicolumn{8}{l}{\textit{CTLE — learned codebook (this work)}} \\
CTLE + K-means      &  4   &  7.61 & $7.98\times$  & 2.957 & 19.23           & 19{,}239            & \phantom{0}73.6 \\
CTLE + GA$^\dagger$     &  4   &  7.61 & $7.98\times$  & 2.944 & 18.98           & 19{,}114            & 207.7 \\
CTLE + PSO$^\dagger$    &  4   &  7.61 & $7.98\times$  & 2.768 & 15.93           & 16{,}827            & 198.3 \\
\bottomrule
\end{tabular}
%
\smallskip
\begin{minipage}{\textwidth}
  \footnotesize
  \textbf{INT4U collapses} (PPL\,$>$11\,000 / 533\,000) because a single global scale
  cannot represent the heterogeneous distribution of the 9.2\,M-entry embedding matrix.
  \textbf{INT4BW} recovers near-FP32 quality
  (TS: $+14.5\,\%$; WT2: $+16.2\,\%$ vs.\ FP32) at the cost of 25\,\% larger footprint
  (9.51\,MB vs.\ 7.61\,MB) and a 31\,\% throughput penalty on RISC-V
  (see Table~\ref{tab:hw}).
  \textbf{CTLE-PSO} achieves competitive PPL at identical 4-bit storage and
  streams weights directly from Flash with a fixed 64-byte LUT per tensor,
  eliminating per-group scale-table overhead at inference time.
  PSO reduces PPL by \textbf{17.2\,\%} vs.\ K-means on TinyStories and
  \textbf{12.5\,\%} on WikiText-2, confirming the improvement is
  corpus-independent.
\end{minipage}
\end{table*}


\begin{table*}[t]
\centering
\caption{Quantization method comparison on TinyStories-15M.
         All CTLE/INT4 methods use 4-bit weight storage.
         CTLE uses $k{=}16$ codebook entries per tensor;
         INT4U uses one global scale per tensor;
         INT4BW uses one scale per group of $G{=}32$ columns
         ($\approx$5\,bits/weight effective).
         NLL and PPL\,(TS) are evaluated on our 10-sentence TinyStories corpus
         (192 tokens scored).
         PPL\,(WT2) is evaluated on the WikiText-2 test split
         (325,693 tokens, non-overlapping windows of 256 tokens, stride\,=\,256).
         Both corpora use the LLaMA SentencePiece tokenizer.
         \textbf{Important:} TinyStories-15M was trained on children-story text,
         not Wikipedia; WikiText-2 PPL values reflect domain mismatch as well as
         compression error and are \emph{not} directly comparable to numbers
         reported for general-purpose LLMs (OPT, LLaMA).
         $\dagger$~PSO/GA use K-means warm start.}
\label{tab:ablation}
\setlength{\tabcolsep}{5pt}
\renewcommand{\arraystretch}{1.15}
\begin{tabular}{lcccccccc}
\toprule
\textbf{Method} & \textbf{Bits/w} & \textbf{Size (MB)} & \textbf{Ratio} &
\textbf{NLL\,$\downarrow$} & \textbf{PPL (TS)\,$\downarrow$} &
\textbf{PPL (WT2)\,$\downarrow$} & \textbf{Time (s)} \\
\midrule
FP32 (baseline)     & 32   & 60.77 & $1.00\times$  & 2.236 & \phantom{00}9.35  & \phantom{0}6{,}127  & —     \\
\midrule
\multicolumn{8}{l}{\textit{Standard INT4 quantization (baselines)}} \\
INT4 Uniform        &  4   &  7.61 & $7.99\times$  & 9.312 & 11{,}070          & 533{,}069           & \phantom{0}${<}1$ \\
INT4 Block-wise     & $\approx$5 & 9.51 & $6.39\times$ & 2.371 & \textbf{10.71} & \textbf{\phantom{0}7{,}120} & \phantom{0}${<}1$ \\
\midrule
\multicolumn{8}{l}{\textit{CTLE — learned codebook (this work)}} \\
CTLE + K-means      &  4   &  7.61 & $7.98\times$  & 2.957 & 19.23           & 19{,}239            & \phantom{0}73.6 \\
CTLE + GA$^\dagger$     &  4   &  7.61 & $7.98\times$  & 2.944 & 18.98           & 19{,}114            & 207.7 \\
CTLE + PSO$^\dagger$    &  4   &  7.61 & $7.98\times$  & 2.768 & 15.93           & 16{,}827            & 198.3 \\
\bottomrule
\end{tabular}
%
\smallskip
\begin{minipage}{\textwidth}
  \footnotesize
  \textbf{INT4U collapses} (PPL\,$>$11\,000 / 533\,000) because a single global scale
  cannot represent the heterogeneous distribution of the 9.2\,M-entry embedding matrix.
  \textbf{INT4BW} recovers near-FP32 quality
  (TS: $+14.5\,\%$; WT2: $+16.2\,\%$ vs.\ FP32) at the cost of 25\,\% larger footprint
  (9.51\,MB vs.\ 7.61\,MB) and a 31\,\% throughput penalty on RISC-V
  (see Table~\ref{tab:hw}).
  \textbf{CTLE-PSO} achieves competitive PPL at identical 4-bit storage and
  streams weights directly from Flash with a fixed 64-byte LUT per tensor,
  eliminating per-group scale-table overhead at inference time.
  PSO reduces PPL by \textbf{17.2\,\%} vs.\ K-means on TinyStories and
  \textbf{12.5\,\%} on WikiText-2, confirming the improvement is
  corpus-independent.
\end{minipage}
\end{table*}
